% ==========================================================================
%  Chapter 2: History of User-Mode Linux
% ==========================================================================

\chapter{History of User-Mode Linux}
\label{ch:history}

\epigraph{``I'm doing a (free) operating system (just a hobby, won't be big and
professional like gnu) for 386(486) AT clones.''}{---Linus Torvalds, 1991~\cite{torvalds1991}}

\noindent
User-Mode Linux did not emerge in a vacuum.  Its existence is a consequence of
design decisions made years before its conception---most notably, the Linux
kernel's clean separation of architecture-dependent and architecture-independent
code.  This chapter traces the full arc: from the early Linux kernel and the
debugging problems that motivated \UML{}, through its creation and golden era, into
a decade of decline, and finally to the revival that produced the modern
multi-backend architecture described in this report.

% ------------------------------------------------------------------
\section{The Linux Kernel Before UML (1991--1999)}
\label{sec:history-prelude}
% ------------------------------------------------------------------

On October 5, 1991, Linus Torvalds posted the famous announcement to
\texttt{comp.os.minix} describing his ``free minix-like kernel sources for
386-AT''~\cite{torvalds1991}.  Linux 0.01 was a monolithic kernel that ran
exclusively on the Intel 80386 processor.  It supported a single address space
layout, a single disk controller, and a handful of system calls.  But it
compiled, it booted, and it was free.

Over the next eight years, Linux grew at an extraordinary rate.  The 1.0 release
(March 1994) was the first declared stable; 2.0 (June 1996) added symmetric
multiprocessing (SMP) support and multiple architecture ports---Alpha, SPARC,
MIPS, m68k, PowerPC.  The 2.2 release (January 1999) brought improved SMP
scalability, better networking, and the beginnings of enterprise-grade
reliability.  By the end of the decade, Linux was running on everything from
embedded routers to IBM mainframes, and kernel development involved hundreds of
contributors worldwide.

\subsection{The \texttt{arch/} Abstraction}
\label{sec:history-arch-abstraction}

The most architecturally consequential decision of this era was the directory
structure that emerged to support multiple hardware platforms.  Rather than
scattering \texttt{\#ifdef} conditionals throughout the code, Linux adopted a
clean split: the \umlpath{arch/} directory contained all architecture-specific
code---trap handling, context switching, page table manipulation, signal
delivery, system call entry---while the rest of the tree
(\umlpath{kernel/}, \umlpath{mm/}, \umlpath{fs/}, \umlpath{net/},
\umlpath{drivers/}) remained architecture-independent.

Each architecture port implemented a well-defined interface: functions like
\texttt{switch\_to()}, \texttt{copy\_thread()}, \texttt{start\_thread()}, and
page table accessors (\texttt{pgd\_val()}, \texttt{pte\_val()}, etc.).  The
interface was never formally specified---it evolved organically as new
architectures were added---but it was remarkably clean.  A new port meant
implementing perhaps 50--100 key functions and data structures, after which the
entire kernel (scheduler, VFS, networking stack, memory manager) would work
without modification.

This abstraction was conceived to support hardware architectures: x86, Alpha,
SPARC, ARM.  No one anticipated that someone would write an ``architecture'' where
the hardware was another copy of Linux itself.

\subsection{The Debugging Problem}
\label{sec:history-debugging}

Kernel debugging in the 1990s was, to put it charitably, primitive.  The primary
tools were:

\begin{description}[leftmargin=2em,labelindent=1em]
  \item[\texttt{printk()}] The kernel's printf equivalent.  Developers added
    \texttt{printk()} calls, rebuilt the kernel, rebooted, reproduced the bug,
    read the console output.  This cycle could take minutes to hours depending on
    the bug.

  \item[Serial console] A step up: connect two machines with a serial cable, use
    one to observe the other's console output.  Essential for debugging panics
    where the screen was already dead.

  \item[\texttt{kgdb}] An external patch (never merged into mainline during this
    era) that allowed GDB to connect to a running kernel via serial line.  It
    worked, but it required a second machine, serial hardware, and careful setup.
    Breakpoints in interrupt handlers were hazardous.

  \item[\texttt{kdb}] SGI's kernel debugger, also out-of-tree.  More capable
    than \texttt{kgdb} for some workflows, but equally dependent on specialized
    hardware setups.

  \item[Crash analysis] When the kernel panicked, you got a register dump and a
    call trace.  Interpreting these required intimate knowledge of the
    architecture's calling conventions and the ability to cross-reference
    addresses against \texttt{System.map}.
\end{description}

\noindent
What you could \emph{not} do was the thing that every application programmer
took for granted: run the code under a debugger.  You could not set a breakpoint
in the scheduler and single-step through a context switch.  You could not attach
\texttt{valgrind} to the memory manager.  You could not run the kernel under
\texttt{strace} to see what system calls it made.  The kernel was the bottom of
the stack---there was nothing underneath it to provide these services.

This asymmetry between kernel and userspace development productivity was a
constant source of frustration.  It also meant that certain classes of bugs---race
conditions, memory corruption, subtle locking errors---were disproportionately
expensive to diagnose and fix.  The kernel needed a way to be a \emph{process}.

% ------------------------------------------------------------------
\section{Jeff Dike and the Creation of UML (1999--2001)}
\label{sec:history-creation}
% ------------------------------------------------------------------

Jeff Dike was a kernel developer at VA Linux Systems in the late 1990s.  VA
Linux was one of the first companies to build a business around Linux, and its
engineers were deeply involved in kernel development.  Dike had been thinking
about the debugging problem and, more broadly, about how to make the kernel more
accessible for development and experimentation.

His key insight was deceptively simple: \emph{the Linux kernel's
\texttt{arch/} abstraction is clean enough that you can write an
``architecture'' where the ``hardware'' is the host kernel's system call
interface.}

Consider what an architecture port must provide: a way to create and switch
between address spaces, a way to handle interrupts and traps, a way to start and
stop the CPU, a way to manage memory pages.  All of these have direct analogs in
the POSIX API:

\begin{itemize}[nosep]
  \item Address spaces $\to$ \syscall{mmap}, \syscall{mprotect}, \syscall{munmap}
  \item Interrupts $\to$ POSIX signals (SIGALRM for the timer tick, SIGIO for I/O)
  \item Context switching $\to$ \syscall{setjmp}/\syscall{longjmp} or signal stack manipulation
  \item Trap handling $\to$ \syscall{ptrace} to intercept guest system calls
  \item Physical memory $\to$ a large \syscall{mmap}'d file (the ``physical memory'' file)
  \item Block devices $\to$ host files accessed via \syscall{read}/\syscall{write}
  \item Network devices $\to$ host TAP/TUN interfaces or Unix domain sockets
\end{itemize}

\noindent
The resulting system would not be an emulator.  It would not translate
instructions or simulate hardware.  It \emph{would be} the Linux kernel---the
same source code, compiled for a different ``architecture.''  The kernel's
scheduler would schedule processes; its VFS would manage file systems; its TCP/IP
stack would route packets.  The only difference was that instead of running on
bare metal (or in a hypervisor), it ran as a collection of ordinary Linux
processes.

\subsection{The Initial Announcement}
\label{sec:history-announcement}

Dike announced User-Mode Linux on the Linux Kernel Mailing List (LKML) in
1999.  The initial response was a mixture of fascination and skepticism.  Some
developers immediately grasped the implications for debugging and testing.
Others questioned whether the performance overhead would make it practical.
A few doubted whether the \texttt{arch/} interface was really clean enough to
support such an unconventional ``architecture.''

Dike persisted.  By 2000, he had a working implementation that could boot a
full Linux system, run a shell, and execute userspace programs.  He presented
the work at the 4th Annual Linux Showcase and Conference (ALS) in
Atlanta~\cite{dike2000}.

\subsection{The Tracing Thread (TT) Mode}
\label{sec:history-tt}

The first \UML{} execution engine used a model called ``tracing thread'' (tt)
mode.  The architecture was as follows:

\begin{enumerate}[nosep]
  \item The \UML{} kernel ran as a single host process (the ``tracing thread'').
  \item Each guest userspace process was a separate host process.
  \item The tracing thread used \syscall{ptrace} to attach to every guest
    process as a debugger.
  \item When a guest process made a system call, the host kernel delivered a
    \texttt{SIGTRAP} to the tracing thread (via the ptrace notification
    mechanism).
  \item The tracing thread intercepted the system call, extracted its arguments
    from the guest's registers, dispatched it to the \UML{} kernel's system call
    handler, wrote the result back into the guest's registers, and resumed the
    guest.
\end{enumerate}

\noindent
This worked, and it was conceptually elegant: ptrace was designed exactly for
this kind of process inspection and control.  But it had significant drawbacks:

\begin{description}[leftmargin=2em,labelindent=1em]
  \item[Host visibility] Every guest process appeared as a real process on the
    host.  If the guest ran 500 processes, the host's process table showed 500+
    entries.  This was confusing for administrators and made it impossible to
    cleanly distinguish guest and host processes.

  \item[Security] Because guest processes were real host processes, they had
    real PIDs, real signal delivery, and real access to \texttt{/proc/self/}.
    A sufficiently determined attacker inside the guest could potentially
    escape---for example, by exploiting ptrace race conditions or by directly
    invoking host system calls that the tracing thread failed to intercept.

  \item[Performance] Every guest system call required \emph{at minimum} four
    context switches: guest $\to$ host kernel (for the syscall trap) $\to$
    tracing thread (SIGTRAP delivery) $\to$ host kernel (to resume guest) $\to$
    guest.  In practice, the overhead was worse because the tracing thread also
    needed to read and write guest registers via \texttt{PTRACE\_PEEKUSR} and
    \texttt{PTRACE\_POKEUSR}.  Each such operation was itself a system call.

  \item[Signal handling] Delivering a signal to a guest process required
    elaborate coordination between the tracing thread, the host kernel's signal
    machinery, and the guest's signal handling code.  Getting this right was one
    of the most complex parts of the implementation.
\end{description}

\noindent
Despite these limitations, tt mode was a remarkable achievement.  For the first
time, a developer could run the Linux kernel under GDB, set breakpoints in the
scheduler, single-step through context switches, and use \texttt{valgrind} to
detect memory errors.  The productivity impact was immediate and
substantial~\cite{dike2001}.

\subsection{Key Insight: UML Is Not an Emulator}
\label{sec:history-not-emulator}

It is worth emphasizing a point that was frequently misunderstood (and still
is): \UML{} is not an emulator.  Emulators like Bochs or (later) QEMU translate
machine instructions from one ISA to another, or simulate hardware devices.
\UML{} does neither.  The guest kernel code is \emph{compiled natively} for the
host architecture.  When a function like \texttt{schedule()} runs inside \UML{},
it is executing the same compiled x86-64 instructions that it would on bare
metal.  The architecture layer merely provides a different implementation of the
hardware abstraction---one that talks to the host kernel instead of to physical
hardware.

This distinction has profound implications for fidelity.  Code that works in
\UML{} will work on bare metal (modulo hardware-specific drivers), and bugs found
in \UML{} are real kernel bugs.  \UML{} is the kernel, running in a different
environment.

% ------------------------------------------------------------------
\section{The SKAS Evolution (2002--2006)}
\label{sec:history-skas}
% ------------------------------------------------------------------

The limitations of tt mode---particularly the security and host-visibility
problems---motivated a fundamental redesign of the execution engine.  The
result was SKAS: \textbf{S}eparate \textbf{K}ernel \textbf{A}ddress
\textbf{S}pace mode.

\subsection{The Problem with Shared Address Spaces}
\label{sec:history-skas-problem}

In tt mode, the \UML{} kernel code ran in the same address space as the
guest process it was currently serving.  This meant that a guest process
could, in principle, read or corrupt kernel data structures.  The tracing
thread did its best to prevent this, but the fundamental architecture made
strong isolation impossible.

The goal of SKAS was to give the \UML{} kernel its own address space,
separate from all guest processes.  When a guest process made a system call,
execution would switch to the kernel address space, handle the call, and
switch back---just as a real kernel switches between kernel and user mode
via privilege ring transitions.

\subsection{SKAS0: No Host Kernel Patches Required}
\label{sec:history-skas0}

The first SKAS implementation, known as SKAS0, achieved address space
separation using only standard host kernel facilities.  The technique was
ingenious:

\begin{enumerate}[nosep]
  \item The \UML{} kernel ran in a dedicated host process with its own
    address space.
  \item Guest processes each had their own host process (as in tt mode),
    but the tracing thread could now switch the guest process's memory
    mappings to point at the kernel's address space when needed.
  \item Address space switching was accomplished via careful use of
    \texttt{PTRACE\_SYSCALL} to intercept \syscall{mmap} calls and remap
    the process's memory.
\end{enumerate}

\noindent
SKAS0 was a major improvement over tt mode.  Guest processes could no longer
directly access kernel memory, and the security model was substantially
stronger.  The performance cost of the address-space switching was nontrivial,
but the security benefits were considered worth it.

\subsection{SKAS3: The /proc/mm Approach}
\label{sec:history-skas3}

SKAS0's address-space manipulation was clever but slow, requiring multiple
system calls per switch.  Dike proposed a host kernel extension---a
\texttt{/proc/mm} interface---that would allow a tracing process to directly
manipulate another process's address space with a single system call.  This
variant was called SKAS3 (the number referring to the host kernel patch
revision that introduced the necessary support).

With \texttt{/proc/mm}, the tracing thread could:
\begin{itemize}[nosep]
  \item Create a new memory management context via \texttt{/proc/mm}
  \item Map and unmap pages in that context with a single \syscall{ioctl}
  \item Switch a traced process between contexts efficiently
\end{itemize}

\noindent
The performance improvement over SKAS0 was significant---roughly 2x for
system-call-intensive workloads.  However, the approach required running a
patched host kernel, which limited adoption.  Most users preferred the
slightly slower SKAS0 over the burden of maintaining a custom host kernel.

\subsection{SKAS4 and the OLS Talks}
\label{sec:history-skas4}

Dike continued refining the approach.  SKAS4 proposed further optimizations
to the host kernel interface, aiming to minimize the number of context switches
required for guest system call handling.  He presented this work at the Ottawa
Linux Symposium in 2006~\cite{dike-ols2006}, building on his earlier OLS 2003
talk~\cite{dike-ols2003} that had focused on making \UML{} ``ready for the
mainstream.''

The SKAS evolution demonstrated a recurring tension in \UML{}'s history:
achieving good performance often required host kernel modifications, but
requiring host kernel modifications limited adoption.  This tension would
recur with every subsequent execution engine, including the seccomp backend
and the \KVMvTwo{} backend described in this report.

\subsection{The Prentice Hall Book}
\label{sec:history-book}

In 2006, Dike published \emph{User Mode Linux} with Prentice
Hall~\cite{dike2006book}.  The book provided the first comprehensive treatment
of \UML{}'s architecture, covering the tt and SKAS execution models, the memory
system, device drivers, networking, and debugging techniques.  It also included
practical guidance on using \UML{} for development, testing, and virtual
hosting.

The book remains the most detailed single source on \UML{}'s original
architecture, although much of the implementation detail it describes has
since been superseded by the seccomp and \KVMvTwo{} backends.

% ------------------------------------------------------------------
\section{Mainline Merge and the Golden Era (2002--2010)}
\label{sec:history-golden}
% ------------------------------------------------------------------

\UML{} was merged into the Linux 2.5 development tree in 2002 and was
included in the 2.6.0 release (December 2003).  The merge was a significant
event: \UML{} was the first virtualization technology to be included in the
mainline Linux kernel.  Its presence in \umlpath{arch/um/} gave it
legitimacy and ensured that kernel API changes would (in theory) be
propagated to the \UML{} port along with every other architecture.

The years from 2002 to roughly 2010 represented \UML{}'s golden era.  During
this period, it was actively developed, widely used, and considered a key
part of the kernel developer's toolkit.

\subsection{Kernel Development and Debugging}
\label{sec:history-kernel-dev}

The primary use case was exactly what Dike had envisioned: making kernel
development easier.  With \UML{}, a developer could:

\begin{itemize}[nosep]
  \item Boot a full Linux system in seconds, without hardware or a heavy VM
  \item Run the kernel under GDB with full source-level debugging
  \item Use \texttt{valgrind} to detect memory errors in kernel code
  \item Test file system changes without risking data on real disks
  \item Reproduce and debug race conditions in the scheduler
  \item Develop and test networking code with virtual networks
  \item Run kernel unit tests without dedicated hardware
\end{itemize}

\noindent
The fast boot time was particularly valuable.  A \UML{} instance could go from
cold start to a shell prompt in under two seconds on contemporary hardware,
compared to minutes for a full hardware boot.  For the edit-compile-test cycle,
this speed difference was transformative.

\subsection{Netkit and Network Simulation}
\label{sec:history-netkit}

One of \UML{}'s most successful applications was Netkit, a network emulation
tool developed by Maurizio Pizzonia and Massimo Rimondini at Roma Tre
University~\cite{pizzonia2008}.  Netkit used \UML{} instances as virtual network
nodes, connected by virtual Ethernet links implemented via Unix domain sockets.

Netkit enabled educators and researchers to create complex network topologies on
a single machine.  A student could configure a multi-router network with BGP,
OSPF, DNS, and DHCP---all running real Linux networking code---on a laptop.
Netkit was widely adopted in networking courses at universities across Europe
and became one of \UML{}'s most visible applications.

The related project Marionnet~\cite{marionnet} provided a graphical interface
for constructing virtual network topologies, further lowering the barrier to
entry for network education and experimentation.

\subsection{Security Research and Honeypots}
\label{sec:history-honeypots}

\UML{} found a natural audience in the security research community.  Its
lightweight isolation made it useful for:

\begin{description}[leftmargin=2em,labelindent=1em]
  \item[Honeypots] Running deliberately vulnerable systems to attract and study
    attackers.  \UML{}'s low overhead made it practical to run many instances
    on a single host, and its isolation (especially in SKAS mode) provided
    reasonable confidence that a compromised guest could not easily escape to
    the host.

  \item[Malware analysis] Executing suspicious binaries in a controlled
    environment where every system call could be observed.

  \item[Intrusion detection research] Running experimental IDS systems against
    realistic kernel and network stacks.
\end{description}

\subsection{Virtual Hosting}
\label{sec:history-vhosting}

Before containers and lightweight VMs became mainstream, \UML{} was used by
some hosting providers to offer virtual private servers (VPS).  Each customer
received a \UML{} instance with its own root filesystem, network configuration,
and process space.  The overhead was lower than full hardware emulation (which
at the time meant something like VMware or Bochs), and the isolation model was
stronger than chroot-based approaches.

Companies like Linode (in its early days, before switching to Xen) and various
European hosting providers offered \UML{}-based VPS products.  While the
performance overhead of ptrace-based system call interception made \UML{} slower
than native execution, it was adequate for many web hosting and
light-computation workloads.

% ------------------------------------------------------------------
\section{Decline and Maintenance Mode (2010--2020)}
\label{sec:history-decline}
% ------------------------------------------------------------------

\UML{}'s decline was not the result of any single event.  It was the
cumulative effect of multiple technological shifts that eroded its competitive
advantages while \UML{} itself stagnated.

\subsection{Hardware Virtualization Changes the Landscape}
\label{sec:history-hwvirt}

The most consequential shift was the arrival of hardware virtualization
support in commodity x86 processors.  Intel released VT-x (originally
codenamed ``Vanderpool'') in late 2005 on the Pentium 4 662 and
672~\cite{intel-vt}.  AMD followed in 2006 with AMD-V (SVM, originally
``Pacifica'')~\cite{amd-svm}.  These extensions added a new privilege
level---VMX root and non-root modes (Intel) or host and guest modes
(AMD)---that allowed a hypervisor to run unmodified guest operating systems
with near-native performance.

The impact on the virtualization landscape was seismic.  Before VT-x/AMD-V,
running an unmodified guest OS on x86 required either full emulation
(slow), binary translation (complex, patented by VMware), or
paravirtualization (required guest OS modifications, as with Xen).
Hardware virtualization made it possible to run unmodified guests with
overhead measured in single-digit percentages rather than factors of 2--10x.

For \UML{}, this was an existential threat.  \UML{}'s ptrace-based system
call interception imposed an overhead of roughly 1--5\,\textmu{}s per guest
system call.  For a workload making 100,000 system calls per second, this
translated to 0.1--0.5 seconds of overhead per second of execution---a
10--50\% performance penalty on a system-call-intensive workload.  Hardware
virtualization reduced the equivalent overhead to tens of nanoseconds per
VM exit, roughly two orders of magnitude faster.

\subsection{KVM: Virtualization in the Kernel}
\label{sec:history-kvm}

The decisive blow came in February 2007, when Avi Kivity's KVM (Kernel-based
Virtual Machine) was merged into Linux 2.6.20~\cite{kivity2007}.  \KVM{}
leveraged VT-x/AMD-V to turn the Linux kernel itself into a hypervisor.  It
was elegant, minimal (the initial patch was under 10,000 lines), and
\emph{fast}.

\KVM{}'s architecture was the mirror image of \UML{}'s.  Where \UML{} ran a
kernel as a userspace process, \KVM{} ran guest kernels in hardware-supported
VM mode, with the host kernel as the hypervisor.  Where \UML{} used ptrace
to intercept system calls, \KVM{} used VT-x/AMD-V to trap only privileged
operations.  Where \UML{} shared the host's physical memory via \syscall{mmap},
\KVM{} used hardware page tables (EPT/NPT) to provide isolated guest physical
address spaces.

Combined with QEMU~\cite{bellard2005} for device emulation, \KVM{} quickly
became the default virtualization solution for Linux.  It offered near-native
performance for CPU-bound workloads and excellent I/O performance via virtio
paravirtualized drivers.  For the use cases where \UML{} had been the best
available option---kernel development, testing, lightweight VMs---\KVM{}/QEMU
was now competitive and often superior.

\subsection{Containers Capture the Lightweight Market}
\label{sec:history-containers}

The other major technological shift was the rise of OS-level virtualization
(containers).  The Linux kernel's cgroup and namespace infrastructure matured
between 2006 and 2013~\cite{soltesz2007,namespaces}, providing the building
blocks for lightweight isolation: PID namespaces (2008), network namespaces
(2009), user namespaces (2012), and cgroup v1 for resource accounting and
limits.

Docker, released in 2013, packaged these primitives into an accessible
developer tool and ignited the container revolution.  For the use case of
``run an isolated Linux environment with low overhead,'' containers offered
effectively zero additional syscall overhead (the guest ran natively, sharing
the host kernel) and millisecond startup times.

Containers did not provide the kernel-level isolation that \UML{} offered---a
container shared the host kernel, so kernel bugs affected all containers---but
for the vast majority of development and deployment workloads, the trade-off
was acceptable.  The combination of near-zero overhead and excellent tooling
made containers the dominant lightweight isolation mechanism.

\subsection{Code Rot and Declining Maintenance}
\label{sec:history-rot}

As users migrated to \KVM{} and containers, developer attention followed.
\UML{}'s maintenance situation deteriorated:

\begin{itemize}[nosep]
  \item The tt mode was eventually removed, simplifying the codebase but
    also removing a working (if slow) execution engine.
  \item SMP support was never completed; \UML{} remained UP-only for over
    a decade.
  \item Various kernel API changes were not properly propagated to
    \umlpath{arch/um/}, causing build breakages that sometimes persisted
    for entire release cycles.
  \item The SKAS0 mode accumulated technical debt, with workarounds and
    hacks layered on top of the original ptrace-based design.
  \item Test coverage was minimal---there were no automated selftests, and
    manual testing was sporadic at best.
  \item Documentation drifted out of date with the actual implementation.
\end{itemize}

\noindent
Richard Weinberger took over as nominal maintainer, and his efforts kept
\UML{} from being removed entirely.  But the subsystem was in maintenance
mode: patches were mostly limited to fixing build breakages caused by core
kernel API changes, with little new feature development.

The perception within the kernel community reflected this state.  \UML{} was
widely viewed as a legacy subsystem---useful in its time, but superseded by
better alternatives.  Proposals to remove it from the tree surfaced
periodically, though they never gained enough traction to reach consensus.

% ------------------------------------------------------------------
\section{The Linux Virtualization Timeline}
\label{sec:history-virt-timeline}
% ------------------------------------------------------------------

To understand \UML{}'s trajectory, it helps to place it in the broader context
of Linux virtualization history.  The following milestones shaped the landscape
in which \UML{} operated:

\subsection{Xen and Paravirtualization (2003)}
\label{sec:history-xen}

Xen, developed at the University of Cambridge, introduced efficient
paravirtualization for x86~\cite{barham2003}.  By modifying the guest OS to
make ``hypercalls'' instead of privileged instructions, Xen achieved
virtualization overhead as low as 2--5\% for many workloads.  Xen was widely
adopted in cloud computing (notably by Amazon Web Services in its early
infrastructure) and demonstrated that high-performance virtualization was
possible even without hardware support.

Xen's success illustrated a design philosophy opposite to \UML{}'s: rather
than making the kernel run as a userspace process, Xen made the hardware (or
a thin software layer mimicking hardware) available to multiple kernel
instances.  Both approaches achieved isolation, but Xen's approach was more
natural for production deployment.

\subsection{Intel VT-x and AMD-V (2005--2006)}
\label{sec:history-hwext}

Intel's VT-x~\cite{intel-vt} and AMD's SVM~\cite{amd-svm} extensions added
hardware support for virtualization to commodity x86 processors.  The key
features included:

\begin{itemize}[nosep]
  \item A new CPU mode (VMX non-root / guest mode) for running guest code
  \item Automatic trapping of privileged operations to the hypervisor
    (VM exits)
  \item Hardware-managed guest/host state save/restore (VMCS/VMCB)
  \item Extended Page Tables (EPT, Intel, 2008) and Nested Page Tables
    (NPT, AMD, 2007) for efficient guest address translation
\end{itemize}

\noindent
These extensions eliminated the need for binary translation or
paravirtualization, enabling hypervisors to run unmodified guest operating
systems with near-native performance.  They also provided a hardware-enforced
security boundary between guest and host that was far stronger than any
software-based approach, including \UML{}'s ptrace-based isolation.

\subsection{KVM Enters Mainline (2007)}
\label{sec:history-kvm-merge}

\KVM{}'s merge into Linux 2.6.20~\cite{kivity2007} was a watershed moment.
It established that the Linux kernel itself would be the primary hypervisor
for Linux-based virtualization, rather than an external project like Xen or
a proprietary product like VMware.  \KVM{}'s tight integration with the Linux
kernel meant that it benefited automatically from every kernel improvement:
scheduler optimizations improved guest scheduling, memory management
improvements improved guest memory performance, and so on.

\subsection{Cgroups and Namespaces Mature (2008--2013)}
\label{sec:history-cgroups}

The cgroup and namespace infrastructure that would eventually underpin
containers was developed incrementally between 2006 and 2013~\cite{namespaces}.
Key milestones included:

\begin{itemize}[nosep]
  \item 2006: cgroup v1 infrastructure merged
  \item 2008: PID namespaces
  \item 2009: Network namespaces
  \item 2012: User namespaces (allowing unprivileged container creation)
  \item 2013: Docker 0.1 released, combining cgroups, namespaces, and
    union filesystems into an accessible tool
\end{itemize}

\noindent
Containers~\cite{soltesz2007} did not provide the same isolation guarantees
as virtual machines (a kernel exploit in a container compromises the host),
but they offered effectively zero overhead and subsecond startup.  For
development workflows and many production deployments, this trade-off was
overwhelmingly attractive.

\subsection{gVisor: A Different Approach (2018)}
\label{sec:history-gvisor}

Google's gVisor~\cite{gvisor}, released in 2018, represented a fascinating
convergence with some of \UML{}'s original ideas.  gVisor implements a
user-space kernel (called Sentry) that intercepts guest system calls and
reimplements them in Go.  Like \UML{}, gVisor runs guest applications as
processes on the host.  Unlike \UML{}, gVisor does not use the Linux kernel
source code---it is a clean-room reimplementation of the Linux system call
interface.

gVisor's original platform used ptrace for system call interception (just
as \UML{}'s tt mode did), and later moved to a KVM-based platform for
better performance.  In 2023, gVisor introduced Systrap, a SIGSYS-based
platform that uses seccomp-bpf for system call interception---conceptually
identical to the approach that Benjamin Berg would bring to \UML{}.

The parallel evolution is striking: two projects, decades apart, converging
on the same insight that seccomp-bpf provides a faster mechanism for system
call interception than ptrace.

\subsection{Firecracker: MicroVMs for Serverless (2020)}
\label{sec:history-firecracker}

Amazon Web Services released Firecracker in 2020~\cite{agache2020}, a
purpose-built virtual machine monitor (VMM) for serverless computing.
Firecracker uses \KVM{} as its hypervisor but strips away nearly all of
QEMU's device emulation complexity, providing only a minimal set of
paravirtualized devices (virtio-net, virtio-block, a serial console).  The
result is a VMM that can boot a guest in approximately 125\,ms and
consumes less than 5\,MB of memory per VM.

Firecracker's existence demonstrated that the performance gap between VMs
and containers could be narrowed dramatically with careful engineering.  It
also showed that for production workloads requiring strong isolation, the
kernel community had converged on \KVM{} as the foundation.

% ------------------------------------------------------------------
\section{The Revival (2020--2026)}
\label{sec:history-revival}
% ------------------------------------------------------------------

Against this backdrop, several independent threads of development converged
in the early 2020s to revive \UML{}.  The motivations were varied, but they
shared a common recognition that \UML{}'s unique position---a real kernel
running as a userspace process---provided capabilities that no other
virtualization technology could match.

\subsection{Benjamin Berg and the Seccomp Backend}
\label{sec:history-seccomp}

The most consequential development was Benjamin Berg's work on a seccomp-based
execution backend.  The idea, building on the seccomp-bpf
framework~\cite{seccomp-bpf} that Will Drewry had added to the kernel in 2012,
was to replace ptrace-based system call interception with SIGSYS-based
notification.

The mechanism works as follows:
\begin{enumerate}[nosep]
  \item A seccomp-bpf filter is installed that matches all system calls from
    guest processes.
  \item When a guest process makes a system call, the kernel generates a
    \texttt{SIGSYS} signal instead of executing the system call.
  \item The signal handler (running in the \UML{} kernel context) examines
    the system call, dispatches it to the appropriate \UML{} kernel handler,
    and manipulates the guest's register state to reflect the result.
  \item Execution resumes in the guest process with the system call appearing
    to have completed normally.
\end{enumerate}

\noindent
The performance advantage over ptrace was substantial.  Ptrace-based
interception required a minimum of two context switches per system call
(guest $\to$ tracer $\to$ guest) plus additional system calls for register
access.  The seccomp approach handled everything within a single signal
delivery, which the kernel could process largely in-place.  Benchmarks showed
roughly 3--4x improvement in system call throughput.

Berg also redesigned the \UML{} process model around futex-based
synchronization, replacing the signal-heavy IPC that the old ptrace backend
had used.  The result was a cleaner, faster, and more maintainable execution
engine.

The seccomp backend was merged into Linux 6.16 (released 2025), marking the
first major new feature to land in \umlpath{arch/um/} in over a decade.

\subsection{Tiwei Bie and SMP Support}
\label{sec:history-smp}

Tiwei Bie (Intel) developed a series of patches to bring SMP support to
\UML{}.  The v4 series implemented multi-CPU scheduling using host threads
as virtual CPUs, with appropriate locking and synchronization to ensure
correct concurrent execution.

SMP support was a long-standing gap in \UML{}'s capabilities.  Modern
kernel code is thoroughly concurrent, with per-CPU data, RCU, and lock-free
algorithms that behave differently on UP and SMP systems.  Running \UML{}
in UP mode meant that an entire class of concurrency bugs---the class that
is most difficult to diagnose and fix---was invisible.

\subsection{Anton Ivanov and Virtio-UML}
\label{sec:history-virtio}

Anton Ivanov contributed virtio-uml, a transport layer that allows \UML{}
to use virtio devices via vhost-user or direct kernel virtio backends.
This work was significant because it connected \UML{} to the broader virtio
ecosystem: instead of implementing custom device models for each device
type (block, network, console), \UML{} could use the same paravirtualized
device stack that \KVM{}/QEMU, Firecracker, and other VMMs used.

\subsection{Hajime Tazaki and the LKL-into-UML RFC}
\label{sec:history-lkl}

Hajime Tazaki submitted an RFC series (reaching v8) proposing to unify the
Linux Kernel Library (LKL)~\cite{lkl,tazaki2019} into \UML{}.  LKL
compiled the Linux kernel as a library that could be linked into userspace
applications, providing access to the kernel's networking stack, file
systems, and other subsystems without running a full VM.

The LKL-into-\UML{} RFC stalled due to the complexity of reconciling LKL's
library model with \UML{}'s process model, but it demonstrated continuing
interest in the fundamental idea of running kernel code in userspace.  It
also highlighted the need for a cleaner backend abstraction---a need that
the modern \texttt{backend\_ops} architecture addresses.

\subsection{The Coding Assistants Policy}
\label{sec:history-ai}

In April 2026, the Linux kernel documentation was updated with
\texttt{coding-assistants.rst}~\cite{coding-assistants}, establishing
official guidelines for AI-assisted kernel development.  While not specific
to \UML{}, this policy change facilitated the rapid development of
infrastructure code (selftests, documentation, build tooling) that had been
a persistent bottleneck.

\subsection{The v2 Architecture}
\label{sec:history-v2}

The convergence of these efforts---seccomp backend, SMP, virtio, and renewed
community interest---created the conditions for a more fundamental redesign.
The v2 architecture, detailed in Chapter~\ref{ch:v2redesign} of this report,
introduces a \texttt{backend\_ops} abstraction layer that cleanly separates
the \UML{} kernel from its execution engine.  This allows multiple backends
(seccomp, \KVMvTwo{}, and potentially others) to coexist, selected at boot
time, while sharing the rest of the kernel infrastructure.

The \KVMvTwo{} backend, described in Chapter~\ref{ch:kvm}, leverages
the same \KVM{} hardware virtualization that powers production cloud
infrastructure, but uses it to accelerate \UML{} rather than to replace it.
This closes the performance gap with bare-metal execution while preserving
\UML{}'s defining advantage: the kernel runs as a debuggable, instrumentable
userspace process.

% ------------------------------------------------------------------
\section{Timeline}
\label{sec:history-timeline-fig}
% ------------------------------------------------------------------

Figure~\ref{fig:timeline} presents a visual timeline of the major events
described in this chapter, placing \UML{}'s development in the context of
the broader Linux virtualization ecosystem.

\begin{figure}[htbp]
\centering
\begin{tikzpicture}[
  x=0.32cm,
  event/.style={fill=UMLBlue!8, draw=UMLBlue!40, rounded corners=2pt,
                font=\scriptsize, text width=3.8cm, align=left, inner sep=3pt},
  virt/.style={fill=UMLTeal!8, draw=UMLTeal!40, rounded corners=2pt,
               font=\scriptsize, text width=3.8cm, align=left, inner sep=3pt},
  uml/.style={fill=UMLGreen!8, draw=UMLGreen!50, rounded corners=2pt,
              font=\scriptsize, text width=3.8cm, align=left, inner sep=3pt},
  decline/.style={fill=UMLRed!6, draw=UMLRed!30, rounded corners=2pt,
                  font=\scriptsize, text width=3.8cm, align=left, inner sep=3pt},
  yearmark/.style={font=\tiny\bfseries, color=UMLMuted},
  connector/.style={thin, color=UMLMuted!60},
]

% --- Axis ---
\draw[thick, color=UMLInk!50, -{Stealth[length=5pt]}] (0,0) -- (36,0);

% Year marks
\foreach \y/\label in {%
  0/1991, 2.6/1993, 5.1/1995, 7.7/1997, 10.3/1999,
  12.9/2001, 15.4/2003, 17.9/2005, 20.6/2007, 23.1/2009,
  25.7/2011, 28.3/2013, 30.9/2015, 33.4/2018, 34.9/2020,
  35.4/2022, 35.7/2024, 36/2026} {
  \draw[color=UMLMuted!40] (\y, -0.15) -- (\y, 0.15);
  \node[yearmark, below] at (\y, -0.2) {\label};
}

% --- Upper events (UML and kernel) ---

% Linux 0.01
\node[event, anchor=south west] (linux) at (0, 1.8)
  {\textbf{1991} Linux 0.01 released};
\draw[connector] (0, 0.1) -- (0, 1.8);

% Linux 2.0 SMP
\node[event, anchor=south west] (smp) at (4.5, 3.5)
  {\textbf{1996} Linux 2.0: SMP,\\multi-arch};
\draw[connector] (6.4, 0.1) -- (6.4, 3.5);

% UML announced
\node[uml, anchor=south west] (uml-ann) at (9.3, 1.8)
  {\textbf{1999} Jeff Dike announces UML (tt mode)};
\draw[connector] (10.3, 0.1) -- (10.3, 1.8);

% UML ALS 2000
\node[uml, anchor=south west] (als) at (11, 3.8)
  {\textbf{2000} UML presented at ALS};
\draw[connector] (11.6, 0.1) -- (11.6, 3.8);

% UML mainline
\node[uml, anchor=south west] (mainline) at (14.2, 1.8)
  {\textbf{2002--03} UML merged into\\Linux 2.6.x mainline};
\draw[connector] (14.2, 0.1) -- (14.2, 1.8);

% SKAS
\node[uml, anchor=south west] (skas) at (14.2, 5.0)
  {\textbf{2002--06} SKAS0/3/4\\evolution};
\draw[connector] (15.4, 0.1) -- (15.4, 5.0);

% Dike book
\node[uml, anchor=south west] (book) at (17.5, 3.5)
  {\textbf{2006} Dike publishes\\UML book (Prentice Hall)};
\draw[connector] (17.9, 0.1) -- (17.9, 3.5);

% Netkit
\node[uml, anchor=south west] (netkit) at (20.2, 5.0)
  {\textbf{2008} Netkit: UML for\\network education};
\draw[connector] (20.6, 0.1) -- (20.6, 5.0);

% UML decline
\node[decline, anchor=south west] (rot) at (25, 3.5)
  {\textbf{2010--20} Maintenance mode;\\code rot, no SMP};
\draw[connector] (27, 0.1) -- (27, 3.5);

% Seccomp backend
\node[uml, anchor=south west] (seccomp) at (33.5, 5.0)
  {\textbf{2025} Seccomp backend\\merged (Linux 6.16)};
\draw[connector] (34.9, 0.1) -- (34.9, 5.0);

% v2 redesign
\node[uml, anchor=south west] (v2) at (34, 1.8)
  {\textbf{2026} backend\_ops +\\kvm-v2 + SMP};
\draw[connector] (36, 0.1) -- (36, 1.8);

% --- Lower events (competing technologies) ---

% Xen
\node[virt, anchor=north west] (xen) at (14.2, -1.5)
  {\textbf{2003} Xen para-\\virtualization};
\draw[connector] (15.4, -0.1) -- (15.4, -1.5);

% VT-x
\node[virt, anchor=north west] (vtx) at (17, -3.3)
  {\textbf{2005--06} Intel VT-x,\\AMD-V released};
\draw[connector] (17.9, -0.1) -- (17.9, -3.3);

% KVM
\node[virt, anchor=north west] (kvm) at (19.5, -1.5)
  {\textbf{2007} KVM merged into\\Linux mainline};
\draw[connector] (20.6, -0.1) -- (20.6, -1.5);

% Docker
\node[virt, anchor=north west] (docker) at (28, -1.5)
  {\textbf{2013} Docker released;\\containers go mainstream};
\draw[connector] (28.3, -0.1) -- (28.3, -1.5);

% gVisor
\node[virt, anchor=north west] (gvisor-node) at (32.5, -3.3)
  {\textbf{2018} gVisor: user-space\\kernel (Google)};
\draw[connector] (33.4, -0.1) -- (33.4, -3.3);

% Firecracker
\node[virt, anchor=north west] (fc) at (33.5, -1.5)
  {\textbf{2020} Firecracker microVMs};
\draw[connector] (34.9, -0.1) -- (34.9, -1.5);

% --- Legend ---
\node[anchor=north west, font=\scriptsize\bfseries, color=UMLInk] at (0, -5.0) {Legend:};
\node[uml, anchor=west, minimum height=0.4cm, text width=1.8cm] at (2.5, -5.0) {UML event};
\node[virt, anchor=west, minimum height=0.4cm, text width=2.5cm] at (8, -5.0) {Competing technology};
\node[event, anchor=west, minimum height=0.4cm, text width=2.5cm] at (14, -5.0) {Linux kernel milestone};
\node[decline, anchor=west, minimum height=0.4cm, text width=1.8cm] at (20, -5.0) {Decline period};

\end{tikzpicture}

\caption{Timeline of User-Mode Linux development in the context of Linux
  kernel and virtualization history (1991--2026).  Upper events relate to
  \UML{} and the Linux kernel; lower events show competing virtualization
  technologies that shaped \UML{}'s trajectory.}
\label{fig:timeline}
\end{figure}

% ------------------------------------------------------------------
\section{Lessons from History}
\label{sec:history-lessons}
% ------------------------------------------------------------------

The arc of \UML{}'s history offers several lessons relevant to the v2
redesign:

\paragraph{Architecture matters more than performance.}
\UML{}'s original \texttt{arch/} insight---that a well-abstracted interface
can be reimplemented on top of very different substrates---was correct in
1999, correct in 2007 when applied to \KVM{}, and remains correct today.
The v2 \texttt{backend\_ops} abstraction applies the same principle within
\UML{} itself: abstracting the execution engine so that new backends (seccomp,
\KVM{}, and future possibilities) can be added without rewriting the rest of
the architecture.

\paragraph{The ptrace tax was real.}
\UML{}'s decade-long decline was not due to a fundamental flaw in the concept
of running a kernel as a userspace process.  It was due to a specific
implementation bottleneck: ptrace-based system call interception was too slow.
The seccomp backend demonstrates that with a better interception mechanism,
the concept remains viable.  The \KVMvTwo{} backend goes further, eliminating
the interception overhead entirely for most operations.

\paragraph{Isolation without hardware is still valuable.}
Containers have no kernel-level isolation.  Hardware VMs provide strong
isolation but require VT-x/AMD-V and cannot be run inside many restricted
environments (CI/CD containers, nested VMs, some cloud instances).  \UML{}
occupies a unique middle ground: real kernel isolation, implemented purely in
software, runnable anywhere a Linux process can run.  This niche has expanded
rather than contracted as CI/CD pipelines and restricted environments have
become ubiquitous.

\paragraph{Debugging the kernel will always matter.}
No amount of hardware virtualization performance eliminates the need to debug
the kernel.  \KVM{}/QEMU provides a debuggable target via GDB's remote
protocol, but the developer experience is qualitatively different from running
the kernel as a native process.  \UML{} allows kernel code to be debugged with
the same tools, workflows, and muscle memory that developers use for userspace
applications.  This advantage is permanent.

\paragraph{Maintenance requires a constituency.}
\UML{}'s decade of rot was a maintenance problem, not a technical one.  The
technology had no organized user community demanding features, no companies
depending on it for production workloads, and no continuous integration
preventing regressions.  The v2 effort addresses this by integrating kselftests,
providing comprehensive documentation, and connecting \UML{} to use cases
(kernel fuzzing, CI/CD, AI-assisted development) that create ongoing
demand for maintenance.

\bigskip
\noindent
With this historical context established, we turn in the next chapter to the
technical foundations on which \UML{} is built: the Linux process model,
the ptrace and seccomp interfaces, and the \KVM{} API.
